\subsection{Blind-execution instance optimality}
\label{sec:blind}

We finally consider blind execution.  A job that is not tested runs for its
actual processing time, so $r(p)=p$ and the average direct work is the mean
processing time.  The optimal fluid algorithm has four blocks: a
maximum-density testing block, tested medium jobs, jobs run blindly, and
tested long jobs.  We now compute its cost explicitly.

\begin{theorem}[Blind-execution instance optimality]
\label{thm:blind-instance}
Fix $L<\infty$.  There exist randomized nonanticipating algorithms
$\mathcal A^*_{n,L}$, depending only on $n,L$, and a deterministic sequence
$\eps_L(n)\to0$ such that every vector $p\in[0,L]^n$ satisfies
\begin{equation}
 \EE\ALG_{\mathcal A^*_{n,L}}(p)
 \le n^2\PhiBE(D[p])+\eps_L(n)n^2.
 \label{eq:intro-blind-upper}
\end{equation}
Conversely, every randomized nonanticipating algorithm $\mathcal A$, even if
it is told the multiset of entries of $p$, satisfies
\begin{equation}
 \EE\ALG_{\mathcal A}(p)
 \ge n^2\PhiBE(D[p])-\eps_L(n)n^2.
 \label{eq:intro-blind-lower}
\end{equation}
The benchmark $\PhiBE$ is defined explicitly below, and one may take
\[
 \eps_L(n)=O_L\!\left(n^{-1/6}\sqrt{\ln(n+2)}\right).
\]
\end{theorem}

\subsubsection{The four-block benchmark}
\label{sec:blind-benchmark}

Let $D$ be a probability distribution on $[0,L]$ and put
\begin{equation}
 \mu=\int p\,dD(p).
 \label{eq:blind-mean}
\end{equation}
Thus blindly executing one unit of well-mixed job mass uses work $\mu$.
Apply \cref{lem:maximum-density-module}, choose its threshold selector $g_*$,
and let $\nu$ be the residual measure from
\cref{eq:common-module-data}.  Write
\[
 a=a_D(g_*),\qquad
 \ell=\int_{(0,L]}p g_*(p)\,dD(p).
\]
The maximum-density identity is
\begin{equation}
 a\tau_D-\ell=1,
 \qquad \tau_D=\frac{1+\ell}{a}.
 \label{eq:blind-test-module}
\end{equation}
Hence a full unit of tests together with the selected short outcomes is a
divisible block of completion mass $a$, work $1+\ell$, and work per
completion $\tau_D$.

For a fixed tested fraction $q$, blind execution is the common fluid model
with alternative work $v=\mu$.  By \cref{lem:test-module-envelope}, its
optimal block collection is
\begin{equation}
 qa\delta_{\tau_D}+q\nu+(1-q)\delta_\mu.
 \label{eq:blind-fluid-blocks}
\end{equation}
If $\tau_D\ge\mu$, the direct block is no slower per completion than any
testing block, so the optimum is pure blind execution and has cost $\mu/2$.
Suppose now that $\tau_D<\mu$.  Sorting the blocks in
\cref{eq:blind-fluid-blocks} by work per completion gives, in order,
\begin{equation}
 \begin{aligned}
 \underbrace{\text{tests and selected }p<\tau_D}_{\text{testing block}}
 &\ \longrightarrow\
 \underbrace{\text{tested }\tau_D<p<\mu}_{\text{SPT order}}\\[-2pt]
 &\ \longrightarrow\
 \underbrace{\text{blindly executed jobs}}_{\text{work }\mu}
 \ \longrightarrow\
 \underbrace{\text{tested }p>\mu}_{\text{SPT order}}.
 \end{aligned}
 \label{eq:blind-four-block-order}
\end{equation}

We next write the area of this completion curve explicitly.  To avoid
inessential boundary notation, first suppose that $D$ has no mass at
$\tau_D$ or $\mu$, and set
\begin{align*}
 \mathcal L&=\{p<\tau_D\},&
 \mathcal M&=\{\tau_D<p<\mu\},&
 \mathcal H&=\{p>\mu\},\\
 a&=D(\mathcal L),&
 \ell&=\int_{\mathcal L}p\,dD(p),&
 \ell_M&=\int_{\mathcal M}p\,dD(p),\\
 d&=D(\mathcal H),&
 D_M&=D|_{\mathcal M},&D_H&=D|_{\mathcal H}.
\end{align*}
For $\bar q=1-q$, remaining-mass accounting gives
\begin{align}
 F_{\mathsf{BE},D}(q)
 &= (1+\ell)\left(q-\frac{aq^2}{2}\right)
    +q\ell_M(\bar q+qd)+q^2\SPT(D_M)\notag\\
 &\quad+\mu\left(\bar q qd+\frac{\bar q^2}{2}\right)
    +q^2\SPT(D_H).
 \label{eq:blind-F}
\end{align}
The five terms describe, respectively, the testing block, the external and
internal costs of the tested medium jobs, the blind block, and the internal
cost of the tested long jobs.  If $D$ has mass at either boundary, split that
mass between the adjacent blocks; the value is unchanged.

Expansion yields
\begin{equation}
 F_{\mathsf{BE},D}(q)=\frac\mu2+Aq+Bq^2,
 \label{eq:blind-quadratic}
\end{equation}
where
\begin{align}
 A&=1+\ell+\ell_M+\mu(d-1)
   =1-\int(\mu-p)^+\,dD(p),
 \label{eq:blind-A}\\
 B&=-\frac{a(1+\ell)}2+\ell_M(d-1)+\SPT(D_M)
   +\mu\left(\frac12-d\right)+\SPT(D_H).
 \label{eq:blind-B}
\end{align}
The feasible set in \cref{eq:common-fluid-envelope} is affine in $q$.
Consequently its completion envelope is concave in $q$, its unfinished-mass
area is convex, and $B\ge0$.  Moreover, when $\tau_D<\mu$,
\cref{eq:common-threshold-equation} gives $A<0$.  A minimizing tested
fraction is therefore
\begin{equation}
 q^*=1\quad(B=0),
 \qquad
 q^*=\min\{1,-A/(2B)\}\quad(B>0).
 \label{eq:blind-qstar}
\end{equation}
We define
\begin{equation}
 \PhiBE(D)=
 \begin{cases}
  \displaystyle\min_{0\le q\le1}F_{\mathsf{BE},D}(q),&\tau_D<\mu,\\[3pt]
  \mu/2,&\tau_D\ge\mu.
 \end{cases}
 \label{eq:blind-Phi}
\end{equation}

\begin{proof}[Proof of \cref{thm:blind-instance}]
Take $r(p)=p$ and $Q=[0,1]$ in
\cref{thm:common-fluid-instance-transfer}.  Then $r$ is one-Lipschitz,
$v_r(D)=\mu$, and \cref{eq:blind-fluid-blocks,eq:blind-Phi} show that
$\Phi_{r,Q}(D)=\PhiBE(D)$.  The two conclusions and
the stated rate now follow from
\cref{eq:common-transfer-upper,eq:common-transfer-lower,eq:common-transfer-rate}.
\end{proof}

\subsubsection{An example requiring all four blocks}

Let
\[
 D(\{0\})=\frac15,\qquad
 D(\{9\})=\frac15,\qquad
 D(\{16\})=\frac35.
\]
Then $\mu=57/5$ and $\tau_D=5$.  Thus tested zeros complete during their
tests, tested jobs of length nine precede the jobs run blindly, and tested
jobs of length sixteen follow them.  Formula \cref{eq:blind-F} becomes
\[
 F_{\mathsf{BE},D}(q)=\frac{57}{10}-\frac{44}{25}q+\frac{11}{10}q^2.
\]
Its minimizer is $q^*=4/5$ and its value is $4.996$.  If all tested nonzero
jobs are instead deferred until after the blind block, the best value is
$5.04$.  The gap is a positive constant at the fluid scale and therefore a
positive constant times $n^2$ for large finite inputs.  In particular, the
simpler order ``testing block, blind block, all tested residuals'' is not
instance optimal.

\subsubsection{Unbounded processing times}

The bound $L$ in \cref{thm:common-fluid-instance-transfer} is necessary.
Without it, blind execution has no finite worst-case ratio.

\begin{theorem}[No finite ratio for unbounded blind execution]
\label{thm:blind-unbounded-impossible}
There are no randomized algorithms $\mathcal A_n$, constant $R<\infty$, and
sequence $\eta(n)\to0$ such that every vector
$p\in[0,\infty)^n$ satisfies
\begin{equation}
 \EE\ALG_{\mathcal A_n}(p)
 \le R\cdot\OPT(p)+\eta(n)n^2.
 \label{eq:blind-unbounded-impossible}
\end{equation}
\end{theorem}

\begin{proof}
Suppose that such a guarantee exists, and run $\mathcal A_n$ on the all-zero
input.  Its offline cost is zero, so if $Z_n$ denotes the expected online
cost, then $Z_n\le\eta(n)n^2=o(n^2)$.  We may assume that the algorithm
finishes every job.

Fix its private seed and order the labels by their first touch.  Every first
touch completes its zero job.  If the job at rank $h$ is tested, that test is
performed while $n-h+1$ jobs are unfinished and contributes $n-h+1$ to the
objective.  Consequently
\begin{equation}
 Z_n\ge\EE\sum_{h:\,\text{rank }h\text{ is tested}}(n-h+1).
 \label{eq:blind-zero-test-area}
\end{equation}

Now choose one label uniformly and change only its length from zero to
$H=n^2$.  Until this exceptional label is first touched, the algorithm sees
the zero-input history.  If its rank $h$ is touched blindly, its execution
occupies the machine for $H$ time units while $n-h+1$ jobs are unfinished.
Averaging over the exceptional label and the private seed gives
\begin{align}
 \EE\ALG_{\mathcal A_n}(p)
 &\ge \frac Hn\,\EE
   \sum_{h:\,\text{rank }h\text{ is blind}}(n-h+1)\notag\\
 &\ge\frac Hn\left(\frac{n(n+1)}2-Z_n\right)
  =H\left(\frac n2-o(n)\right).
 \label{eq:blind-one-long-average}
\end{align}
Some fixed choice of the exceptional label has at least this expected cost.
Its offline optimum runs all zero jobs first and the long job last, and hence
$\OPT(p)=H$.  With $H=n^2$, \cref{eq:blind-one-long-average} contradicts
\cref{eq:blind-unbounded-impossible} for all sufficiently large $n$.
\end{proof}
